%% 第31回編集委員会（物理第5巻）：組 原W8b
%% p5-16.tex  原-16　読みもの　原子の力とつきあう（★2〜4。読みもの。3段にしない：N-1）
%% 中身は計画書 tools/phys5-plan.md の B「原-16」・A8（1回の反応のエネルギーの桁）・A14（原子力の利用とその課題）・A82（崩壊熱）、
%%   E の7（N-7：賛成・反対の立場をとらない。物理と公的資料の事実。被災した人に配慮する。原爆は短く）。
%%   読みもの（例題は「見本の計算」4題。じぶんでとこう・章末問題はなし）。第2巻 p2-16.tex・第4巻 p4-16.tex と同じ形。
%%   中学・高校の公式を本当に導く所がないので「この章で導く中学・高校の公式」の欄は置かない（6-1 の2・K-8）。
%% 節の並び（計画書 B の①〜⑧から変えた所。Z5-9 の照らし合わせ用）：
%%   ①使いみちと始まり（原爆の段落はここ）②使い終わった燃料 ③地層処分（はてな？の答え）④安全のしくみと規制
%%   ⑤福島第一の事故の経過 ⑥地域の暮らしと健康（計画書にない節。原-15 ⑪が「地域の暮らしの今は原-16」と送った中身）
%%   ⑦廃炉と処理水 ⑧放射線の利用 ⑨原子の力とどうつきあうか（計画書の⑧。熱-16 の「社会の選択」と同じ言い方）。
%% 公式は導き直さず、持ち主の章を指した：核分裂 約200 MeV・炭素の 4.1 eV・約5千万倍・1 g で 8.2×10^10 J（原-15 ①・例題1）、
%%   3 GW の炉が1日に約 3 kg（原-15 例題2）、遅発中性子 0.0065・負の反応度係数（原-15 ⑥⑨⑩）、崩壊熱の目安の式
%%   P/P0≈0.066[t^-0.2−(t+T0)^-0.2]（原-15 ⑪ ph:原-15:decayheat。係数・1年＝365.25日も同じ）、ジルコニウムと水蒸気の水素・事故の物理（原-15 ⑪）、
%%   半減期（原-12 ③、3-11④）、ヨウ素131 80日で約1000分の1（原-12 例題2）、年代測定（原-12 ⑦）、放射化（原-13 ⑨）、
%%   非破壊検査・厚さ計（原-13 ⑪）、放射線治療（原-13 ⑫）、Bq・Sv・線量係数・トリチウム 0.000018 μSv/Bq（原-14 ⑤⑧）、
%%   3原則（原-14 ⑦）、自然 2.1 mSv/年・医療 2.6 mSv/年（原-14 ⑩）、100 mSv あたり 0.5 %と原爆の被爆者の調査（原-14 ⑪）、
%%   公衆 1 mSv/年・線量限度は境界ではない（原-14 ⑫）、トリチウムの例題 0.020 mSv（原-14 例題9）、熱-16 の物差しの表（熱-16 ⑨）。
%% 前の巻との約束（計画書 A-6）：Q19（熱-16:332「事故への備えと使い終わった燃料の扱い（原-16）」）→ ②③④⑨（⑨で熱-16 の言い方を回収）。
%%   Q21（熱-16:341「使い終わった燃料の扱いは原-15・原-16」）→ ②③。Q25（熱-16 の道しるべ → 原-16）→ 道しるべの「前に」でなく、⑨と道しるべの「先に」で熱-16 へ返す。
%%   A8（中学：原子力などからエネルギー）→ ①。A14（物理基礎：原子力の利用とその課題）→ 全体。A82（崩壊熱）→ ②・見本の計算1・3。
%% 原-15 から送られた中身：事故の経過（⑤）・廃炉（⑦）・地域の暮らし（⑥）・原爆の短い記述（①）・プールで何年も冷やす理由（②）・
%%   電源と水の確保が設計と規制の中心（④）。原-12 ⑫「保管で放射能を減らす管理（原-16）」→ ②③。
%% 記号（この章）：半減期 T_{1/2}（原-12）、崩壊熱の割合 P/P_0・運転した時間 T_0・止めてからの時間 t（原-15 と同じ）、
%%   実効線量 E（原-14 と同じ。この章ではエネルギーを記号にしない）、水の質量 m・比熱 c・蒸発熱 L（見本の計算3 だけ）。新しい記号カードはなし。
%% 数値はすべて Python で計算した：checks/物理第5巻/原-16/calc_p5_16.py（出力 calc_p5_16_out.txt）、図の位置は fig_p5_16.py、
%%   常用漢字の確かめは kanji_p5_16.py（joyo.txt は2010年の2136字）。
%% 出典（11月11日に作業環境から確かめた。Wikipedia は使っていない）
%%   ・環境省「放射線による健康影響等に関する統一的な基礎資料（令和7年度版）」HTML 版（https://www.env.go.jp/chemi/rhm/current/ の各ページを curl で読んだ）：
%%       01-02-07（ヨウ素131 約8日・セシウム137 約30年）、02-02-01・06-02-04（INES レベル7、チョルノービリと同じ）、
%%       02-02-03（軽水炉の燃料 ウラン235 3〜5 %、ウラン238 からプルトニウム239）、02-02-04（ヨウ素131・セシウム134/137・ストロンチウム90）、
%%       02-02-05（放出量の比較。福島第一は格納容器の大規模な破壊がなかった）、03-04-03（100 mSv あたり 0.5 %は原爆被爆者のデータから）、
%%       03-07-05（原爆被爆者の調査：固形がん、100〜200 mGy より下は不確実）、03-08-04・03-08-05（専門家と住民の対話）、
%%       03-08-18・03-08-19（避難そのものが健康のリスクになり得る。病院・施設の人が避難中・避難後に亡くなった。震災関連死。
%%       OECD/NEA 2024 などの報告）、05-01-02・05-02-07・05-02-08（UNSCEAR 2020年/2021年報告書：「放射線被ばくが直接の原因となるような
%%       将来的な健康影響は検出できそうにない」、甲状腺がんは高感度の超音波検診の結果、避難した人の心血管疾患などは生活の変化による、
%%       福島第一の結果はチョルノービリよりはるかに低度、格納容器が大半を保持）、06-01-01（2011年3月11日 14時46分、M9.0）、
%%       06-01-02〜06-01-04（1〜3号機の交流電源の喪失・燃料の溶融・水素爆発 1号機12日・3号機14日・4号機は3号機から流れこんだ水素）、
%%       06-02-01（19時03分 原子力緊急事態宣言、21時23分 3 km の避難指示、10 km（5万1207人）、12日 15時36分の爆発のあと 20 km）、
%%       06-03-01（2011年12月 冷温停止状態、2019年12月のロードマップ改訂、2024年11月・2025年4月 燃料デブリの試験的取り出し、
%%       2041〜2051年の廃止措置完了の目標）、06-03-02（国・東京電力・規制委員会の体制）、06-03-03（プールの燃料 4号機 2014年12月・3号機 2021年2月、
%%       2024年9月 試験的取り出しに着手）、06-03-04（汚染水 約540 m³/日（2014年5月）→約70 m³/日（2024年度））、
%%       06-03-05〜06-03-07（ALPS：トリチウム以外の62核種、トリチウムは水として存在し除けない、2021年4月13日の基本方針、
%%       2023年8月24日 放出開始、1500 Bq/L 未満・海水で100倍以上にうすめる、告示濃度限度の40分の1、WHO の飲料水の指標の7分の1程度、
%%       告示濃度限度は70歳まで毎日約2 L 飲んで平均 1 mSv/年になる濃さ）、06-03-09（年間 22兆 Bq 未満）、06-03-10（放出口の付近で
%%       30 Bq/L 程度、周辺の海 0.1〜1 Bq/L、WHO 1万 Bq/L、人への影響は 2.1 mSv の約100万分の1〜約7万分の1（2023年2月の評価）、
%%       放出開始後1年の実績で約40万分の1〜約10万分の1）、06-03-11（海域モニタリング、IAEA の支援）、06-03-13（IAEA 包括報告書 2023年7月4日：
%%       国際安全基準に合致、影響は無視できるほど）、09-01-01（除染：取り除く・遮る・遠ざける）、09-02-01・09-02-04（中間貯蔵施設：
%%       大熊町・双葉町、約1600 ha、2045年3月までに県外で最終処分：中間貯蔵・環境安全事業株式会社法）、09-04-02（2011年4月22日 警戒区域・計画的避難区域
%%       20 mSv）、09-05-01・09-05-03・09-05-04（2020年3月までに帰還困難区域を除き解除、特定復興再生拠点区域は 2023年11月までに全て解除、
%%       特定帰還居住区域：2020年代をかけて帰還を望む住民が帰還できるよう）。
%%   ・特定放射性廃棄物の最終処分に関する法律（平成12年法律第117号）第2条第2項（地下 300 m 以上の政令で定める深さの地層に埋設）、
%%       第2条第10項〜第12項・第6条〜第8条（文献調査→概要調査→精密調査）：e-Gov 法令 API https://laws.e-gov.go.jp/api/2/law_data/412AC0000000117 を読んだ。
%%   ・原爆被爆者の寿命調査（約12万人、1950年から）：厚生労働省の審議会資料（放射線影響研究所 寿命調査 第14報の要約、
%%       https://www.mhlw.go.jp/file/05-Shingikai-11201000-Roudoukijunkyoku-Soumuka/0000070377.pdf を pdftotext で読んだ）。
%%   ・原爆の死者の推計（1945年12月末まで）：広島 約14万人（誤差±1万人：広島市の「被爆者 よくある質問」と平和記念資料館の学習ハンドブック）、
%%       長崎 7万3884人（長崎市の資料、長崎原爆資料保存委員会の1950年7月の報告。長崎県の資料・長崎平和推進協会の英語の資料で照合）。
%%       市のサイトは作業環境から届かず、WebSearch の抜粋で照らし合わせた（2026年11月）。紙面は「約14万人」「約7万4千人」。
%%       核分裂の発見（1938年の終わり）・1942年12月2日・1945年に原爆：原-15 の出典（ノーベル賞のサイト）と同じ言い方（「7年もたたない」）。
%%   ・ガラス固化体を30〜50年ほど冷やしてから地層処分、数万年以上生活環境から離す、4万本以上を1か所に：JAEA の ATOMICA の解説、
%%       NUMO の中学生向け教材（2021年）と「地層処分事業の概要」（2025・2026年版）の抜粋（WebSearch で照合。NUMO・JAEA のサイトは届かない）。
%%       京都大学の研究者の資料が「ウラン鉱石なみになるまでの年数」の比べ方を批判していることも見た（紙面は「数万年以上」とだけ書いた）。
%%   ・文献調査：2020年11月に北海道の2町村で始まり、2024年6月に佐賀県の町、2026年5月20日に東京都の南鳥島（全国4例目）、
%%       2026年10月15日に茨城県の市（5例目）：NUMO の広報誌（2026年5月）・プレスリリース（2026年5月20日）、熊本日日新聞（共同通信）の記事、
%%       佐賀新聞の記事の抜粋で照合。処分地は決まっていない。紙面は地名を書かず「2020年から、いくつかの地域で」と書いた（変わりやすいため）。
%%   ・再処理工場（青森県六ヶ所村）は完成していない（2026年の時点。竣工の目標は「2026年度中」だったが審査が遅れている）：
%%       資源エネルギー庁の WG 資料と福井県の資料（2026年6月）、熊本日日新聞（共同通信）の記事の抜粋で照合。
%%   ・使用済燃料の中間貯蔵施設（青森県むつ市、リサイクル燃料貯蔵、乾式貯蔵）は 2024年11月6日に事業開始（国内初）：東京電力の解説ページ、
%%       時事通信の記事、日本原子力産業協会の会長コラム（2024年11月8日）の抜粋で照合。
%%   ・原子力規制委員会の発足 2012年9月19日、新規制基準の施行 2013年7月8日（重大事故への対策を初めて求めた）：東京電力の「新規制基準」の解説、
%%       日本原子力産業協会の資料（2013年7月）、imidas の解説、ATOMICA の要約の抜粋で照合（www.nra.go.jp は届かない）。
%%   ・福島県からの避難者：2012年5月に約16万人（県外約6万人・県内約10万人）、2026年5月1日に 22,625人：福島県「避難者数の推移」、
%%       復興庁の伝承の手引き（2012年5月の 16万4865人）の抜粋で照合。紙面は「約16万人」「約2万3千人」。
%%   ・ウリミバエの根絶（1993年、沖縄県。不妊にした虫を放す方法）：農林水産省 植物防疫所の資料、日本植物防疫協会の資料の抜粋で照合。
%%       不妊にするのに γ線（コバルト60）を使ったことは一般向けの記事で照合した。紙面は「放射線で不妊にした」と書いた。
%%   ・プルトニウム239 の半減期 約2万4千年（24110年）：米国疾病対策センター（CDC）の解説、Jefferson Lab の核種表（WebSearch で照合。
%%       DDEP・NNDC・JAEA は届かない）。計算は 24110年（計画書の 24100年でも 約24万年で同じ）。
%%   ・半減期（IAEA LiveChart の API、ENSDF の基底状態。https://nds.iaea.org/relnsd/v1/data?fields=ground_states&nuclides=137cs ほか。
%%       11月11日に読み、checks/物理第5巻/原-16/iaea_livechart_halflife_20261111.csv に保存）：セシウム137 30.08 年、ヨウ素131 8.0252 日、
%%       プルトニウム239 24110 年（±30）、ストロンチウム90 28.91 年、トリチウム 12.32 年。紙面は環境省の資料と同じ「約30年」「約8日」と、
%%       「約2万4千年」。上の CDC・Jefferson Lab の照合は、この一次資料で確かめ直した。
%%   ・UNSCEAR 2020/2021 報告書が 2019年の終わりまでの科学情報をまとめたこと：www.unscear.org の福島のページ（"up to the end of 2019"）を読んだ。
%%       結論の日本語の言い方は環境省の基礎資料 5.1節・5.2節の訳を使った。
%%   ・崩壊熱が1回の分裂のあと t^-1.2 に比例するとみなせる（Way と Wigner、1948年。Phys. Rev. 73, 1318）：原-15 の出典の Algora ら
%%       arXiv:2007.07918 が Way と Wigner を崩壊熱の最初の見積もりとして引いている。積分で 5[t^-0.2−(t+T0)^-0.2] になることは calc_p5_16.py で
%%       数値積分して確かめた。
%%   ・燃料デブリ 約880トン（推定）、試験的取り出し 2024年11月 約0.7 g・2025年4月 0.187 g（合計 約0.9 g）：nippon.com の記事（東京電力の推定を引く）、
%%       福島県の廃炉の解説、JAEA の分析の資料（2025年）、東京電力の発表（2025年4月24日、約0.2 g）の抜粋で照合。
%%   ・フィンランド（オルキルオト）とスウェーデン（フォルスマルク）が処分の場所を決め処分場を建設：NUMO の広報誌（2022年2月）、原子力環境整備促進・
%%       資金管理センターの資料、Vattenfall の発表（2023年）、共同通信の記事（2026年）の抜粋で照合。
%%   ・オクロの天然原子炉を地層処分のナチュラルアナログとして調べた研究：産業技術総合研究所 地質調査総合センター「地質ニュース」、
%%       日高洋の総説（RADIOISOTOPES 61巻、2012年）、1999年の地球惑星科学連合大会の発表の要旨（CiNii・J-STAGE の目次と要旨）で照合。
%%   ・〔要確認〕なし。〔仮〕（見本の計算で「とします」と書いた値）：3年運転・熱出力 3 GW、プールの水 1000 t・30 ℃、プールの熱は一定。
%% 第32回の確認の指摘（F6）を直した（修5-6。記録は tools/FIX-phys5-review-6.md）。足した出典（2026年11月12日）：
%%   ・タンクの水の約7割（2025年9月の時点）：環境省 基礎資料（令和7年度版）HTML 06-03-06（https://www.env.go.jp/chemi/rhm/current/06-03-06.html を curl で読んだ）。
%%   ・2011年12月の冷温停止：同じ 06-03-01「2011年12月にすべての号機で冷温停止状態を達成」。1〜3号機の冷温停止状態の確認（2011年12月16日）は政府の発表。
%%   ・2号機のプールの燃料の取り出し（2026年6月2日開始、全615体、2028年度中の完了の目標、2026年10月29日で126体）：東京電力の廃炉の進捗資料
%%     （d261029_26-j.pdf）の要約、熊本日日新聞（共同通信）の記事2本、福島県の資料（7月末で28体）を WebSearch で照らし合わせた。
%%     tepco.co.jp は作業環境から届かない（DNS で断られる）。Ver1 の前に一次資料で確かめ直す（H-11）。
%%   ・4〜6号機は定期検査中：東京電力「福島第一原子力発電所の事故の概要」と 2011年3月12日のプレスリリース（WebSearch の抜粋）。
%%   ・避難を続けている人の数（2026年5月1日 22,625人）は、Ver1 の直前に福島県の資料で確かめ直す（H-11。福島県のサイトは作業環境から届かない）。
%%   記号（H-6）：例題3 の温度の上がりは Δθ（T_0・T_{1/2} とぶつけない。原-15 の Δθ と同じ）。本文の「見本の計算N」は囲みの題と同じ「例題N」に。
\makeatletter
% 箱  \txqbox{名前}{(位置)}{はば(mm)}{字}{色}
\newcommand\txqbox[5]{\node[draw=#5!75!black, fill=#5!10, line width=0.8pt, rounded corners=1.4mm, minimum width=#3 mm, minimum height=10mm,
  text width=#3 mm-4mm, align=center, font=\sffamily\footnotesize, inner sep=2pt] (#1) at #2 {#4};}
% 燃料の束（たて長の四角）  \txqfuel{(左下)}
\newcommand\txqfuel[1]{\filldraw[fill=bkOrange!55, draw=bkInk, line width=0.4pt] #1 rectangle ++(2.6,11);}
\makeatother
\tikzset{txq lab/.style={font=\sffamily\footnotesize, text=bkInk, inner sep=1pt},
  txq arr/.style={-{Stealth[length=2.4mm, width=2mm]}, line width=1pt, bkInk!80},
  txq back/.style={-{Stealth[length=2.2mm, width=1.9mm]}, line width=0.8pt, bkBlue!75!black, densely dashed},
  txq dim/.style={bkGray!80!black, line width=0.5pt, {Stealth[length=1.8mm]}-{Stealth[length=1.8mm]}}}

\phchapter[読みもの]{16}{原子の力とつきあう}{使い終わった核燃料は、どうする？}{\bklevelrange{2}{4}}{原子の力とつきあう}
\bkmichishirube{3-11,6-4,熱-8,原-11,原-12,原-13,原-14,原-15}{熱-16,電-16}

\begin{goal}
\begin{itemize}[label={\color{bkTheme}●}, leftmargin=1.6\zw, itemsep=0pt, topsep=1pt]
\item \bklevel{2}\enspace 原子の力の使いみち（発電・医療・工業・農業・研究）と、核分裂の発見から原子炉・原子爆弾・原子力発電までの歴史を説明できる。
\item \bklevel{3}\enspace 使い終わった燃料が熱と放射線を出し続けるわけを、崩壊熱と半減期で説明し、冷やす期間と、人から離しておく期間の桁を見積もれる。
\item \bklevel{3}\enspace 福島第一原子力発電所の事故の経過と、廃炉・処理水のことを、物理の量（崩壊熱・Bq・Sv）と公的な資料の数で説明できる。
\item \bklevel{4}\enspace 放射線や原子力の数を、何を・どの基準で数えた値かに気をつけて読み、物理でわかることと、社会が選ぶことを分けて考えられる。
\end{itemize}
\end{goal}

\begin{hatena}
\noindent\begin{minipage}[c]{0.55\linewidth}\raggedright
原子力発電所では、ウランの燃料を何年か使うと、新しい燃料と取りかえます。使い終わった燃料は、もう電気をつくりません。\par\vspace{3pt}
でも、使い終わった燃料は、すぐに捨てたり、ふつうのごみのように燃やしたりはできません。どうしてでしょう。そして、使い終わった燃料は、このあと、どこへ行くのでしょう。
\end{minipage}\hfill
\begin{minipage}[c]{0.42\linewidth}\centering
\begin{tikzpicture}[x=1mm, y=1mm]
  \fill[bkBlue!14] (2,2) rectangle (50,30);
  \draw[bkInk, line width=1.1pt] (2,34) -- (2,2) -- (50,2) -- (50,34);
  \draw[bkBlue!60, line width=0.6pt] (2,30) -- (50,30);
  \foreach \x in {7,12,17,22,27,32,37,42} {\txqfuel{(\x,4)}}
  \draw[bkInk, line width=0.5pt] (6,3.5) rectangle (45.6,15.5);
  \node[txq lab, anchor=south] at (26,16.5) {使い終わった燃料};
  \node[txq lab, anchor=south] at (26,31) {水のプール};
  \node[font=\sffamily\bfseries\Large, text=phHatena!80!black] at (46,40) {?};
\end{tikzpicture}
\end{minipage}
\end{hatena}

この章は読みものです。じぶんでとこう・章末問題はありません（見本の計算の例題はあります）。これまでの章で、原子核のエネルギー（原-11）、こわれる原子核と半減期（原-12）、放射線の止まり方（原-13）、放射線の量のはかり方（原-14）、原子炉のしくみと崩壊熱（原-15）を学んできました。この章では、その道具を使って、原子の力を人がどう使い、どんな課題とつきあっているかを見ていきます（図\ref{ph:原-16:fig-michi}）。

原子力については、ニュースや本で、いろいろな意見が語られます。この章は、原子力に賛成・反対のどちらかの立場をとることはしません。物理でわかることと、国や国際機関などの公的な資料に書かれている事実を、出典と年をそえて書きます。どう使うかを選ぶのは、社会であり、これからの社会をつくるみなさんです。

この章では、原子爆弾や、原子力発電所の事故のことも取り上げます。この本を読む人の中には、家族や知り合いに被爆した人がいる人や、事故のときに避難した人、事故の影響のあった地域で今も暮らしている人がいると思います。読んでいて、つらいことを思い出したときは、無理をせずに読むのを休んでください。
\begin{bkfigure}
\begin{tikzpicture}[x=1mm, y=1mm]
  \tikzset{src/.style={bk node, text width=22mm, inner xsep=1mm, align=center, minimum height=17mm, font=\sffamily\footnotesize}}
  \node[src] (a) at (0,24) {核分裂\\約200 MeV\\（原-15）};
  \node[src] (b) at (27.2,24) {半減期\\10回で1/1000\\（原-12）};
  \node[src] (c) at (54.4,24) {放射線を\\さえぎる\\（原-13）};
  \node[src] (d) at (81.6,24) {Bq と Sv\\量で考える\\（原-14）};
  \node[src] (e) at (108.8,24) {崩壊熱\\止めても熱\\（原-15）};
  \node[bk node, draw=bkTheme, line width=1pt, text width=96mm, align=left, minimum height=11mm, font=\sffamily\small] (f) at (54.4,-2) {原子の力の使いみちと、課題とのつきあい方（原-16）};
  \foreach \n in {a,b,c,d,e} {\draw[bk link] (\n.south) -- (\n.south |- f.north);}
\end{tikzpicture}
\bkfigcaption{これまでの章の道具が、この章で使われる。}\label{ph:原-16:fig-michi}
\end{bkfigure}

\begin{yattemiyou}[半減期と崩壊熱を、表計算で比べよう]
\begin{enumerate}[label={\bkcircled[bkGreen]{\arabic*}}, leftmargin=2.4\zw, itemsep=1pt, topsep=2pt]
\item 表計算の A 列に 0, 1, 2, …, 20 と、半減期が何回たったかを書く。B 列に $(1/2)^{\text{A 列}}$ を計算する（3-11④）。B 列が 0.001（1000分の1）より小さくなるのは、何回めかを調べる。
\item セシウム137 の半減期を30年として、C 列に「A 列×30」（年）を計算する。1000分の1になるのは、何年後だろう。
\item 別の表に、止めてからの時間 $t$ を 1, 2, 5, 10, 20, 50, 100 年と書き、原-15 ⑪の崩壊熱の目安の式で、3年運転した燃料の熱の割合を計算する（$t$ と運転した時間は秒に直す。1年は \num{31557600} 秒）。1年後から10年後までで、何分の1に減るかを、②の答えと比べよう。
\end{enumerate}
\end{yattemiyou}

\noindent やってみようの答え：半減期の10回分で $(1/2)^{10}=1/1024$、約1000分の1になります（原-12 例題2）。半減期30年のセシウム137 なら、約300年です。一方、崩壊熱は、1年後から10年後までの9年で約7.5分の1にしか減りません。崩壊熱は、半減期の違う何百種類もの分裂片の熱の合計なので、1つの半減期では表せない、ゆっくりした減り方をするのです（原-15 ⑪）。この2つの「時間の物差し」が、使い終わった燃料をどう扱うかを決めています。

\begin{omoidasou}
\bkrecall{原-15}{ウラン235 の核分裂1回で約 \qty{200}{MeV}。炭素が燃えるときの約 \qty{4.1}{eV} の約5千万倍（①）。止めた原子炉でも、分裂片の崩壊熱が $\dfrac{P}{P_0}\approx0.066\,[t^{-0.2}-(t+T_0)^{-0.2}]$（目安、$t$・$T_0$ は秒）で出続ける（⑪）。}
\bkrecall{原-12}{放射性の原子核の数は、半減期 $T_{1/2}$ ごとに半分になる。半減期の10倍で約1000分の1（例題2）。}
\bkrecall{原-14}{実効線量 $E$〔Sv〕は、体への影響を考えるための量。日本の自然放射線は1年に約 \qty{2.1}{mSv}。取りこんだ Bq に線量係数をかけると預託実効線量（⑧）。線量限度は、安全と危険の境界線ではない（⑫）。}
\bkrecall{熱-16}{電気のつくり方には、燃料・二酸化炭素・費用・出したいときに出せるか・資源・安全などの物差しがあり、どれを重く見るかは社会の選択（⑨）。}
\end{omoidasou}

%% ================================================================
\Needspace{16\baselineskip}
\bksection{doguu}{物理の道具}

\bksubsection{1}{原子の力の使いみちと、その始まり}
\leavevmode\label{ph:原-16:sec1}この章で「原子の力」とよぶのは、原子核のもつエネルギーと、原子核から出る放射線のことです。使いみちは、大きく分けて2つあります。1つは、核分裂の熱で電気をつくる\textbf{原子力発電}です。もう1つは、放射線の性質（物を通りぬける・電離する・目印になる）を使う\textbf{放射線の利用}で、病院・工場・農業・研究で使われています（⑧）。\bkidx{げんしりょくはつでん}{原子力発電}

原子力発電の特徴は、エネルギーの桁にあります。原-15 ①で見たように、ウラン235 の核分裂1回で出るエネルギーは約 \qty{200}{MeV} で、炭素の原子1個が燃えるときの約 \qty{4.1}{eV} の約5千万倍です。ウラン235 が \qty{1}{g} 全部分裂すると約 \qty{8.2e10}{J}（原-15 例題1）で、同じ熱を炭素を燃やして得るには、約2.5トンの炭素が要ります。少ない燃料で大きなエネルギーが出ることと、燃やさないので二酸化炭素を出さないことが、原子力発電が使われてきた大きな理由です（熱-16 ⑨）。

\phboxhead{核分裂の発見から、兵器へ}
核分裂が見つかったのは、1938年の終わりのことです（原-15 ①）。1942年12月には、アメリカで、はじめて連鎖反応を制御しながら続ける原子炉が動きました。その一方で、第二次世界大戦の中、連鎖反応を一気に起こす兵器、\textbf{原子爆弾}が作られました。核分裂の発見から7年もたたない1945年8月6日に広島に、8月9日に長崎に、アメリカ軍が原子爆弾を投下しました。爆風と熱線と放射線によって、その年の終わりまでに、広島で約14万人、長崎で約7万4千人が亡くなったと推計されています（広島市・長崎市の推計）。生きのびた人の多くも、けがや病気、その後の生活の苦しみをかかえて生きてきました。\bkidx{げんしばくだん}{原子爆弾}

科学の発見が、わずか数年で、多くの人の命をうばう兵器に結びついたことは、物理を学ぶ人が知っておくべき歴史の事実です。

\phboxhead{被爆した人の調査が、放射線のリスクの土台に}
原爆で被爆した人の健康は、戦後、長い年月にわたって調べられてきました。広島と長崎の約12万人の集団について、1950年から、どんな病気でいつ亡くなったかを追いかける調査（寿命調査）が続けられています。一人一人が受けた線量を推定し、線量とがんの増え方の関係を調べたこの調査は、放射線のリスクを知るための、もっとも大切なデータの1つになっています。原-14 ⑪で学んだ「100 mSv あたり、がんで亡くなる確率が 0.5\,\% 増える」という防護の目安も、このデータをもとにしています（環境省の基礎資料 3.4節・3.7節）。いま私たちが放射線の量を「量で」考えられるのは、調査に長い年月協力してきた被爆者と家族の方々がいるからなのです。

\begin{phmijika}[ことばの使い分け：「被爆」と「被ばく」]
新聞や公的な資料では、原子爆弾の爆撃を受けたことを「被爆」、放射線を受けることを「被ばく」と書き分けることが多くあります。\bkidx{ひばく}{被爆}\bkidx{ひばく}{被ばく}原子爆弾の被害を受けた人は「被爆者」、病院の検査や自然から放射線を受けることは「被ばく」です。この本も、この書き分けにしたがっています。
\end{phmijika}

%% ================================================================
\bksubsection{2}{使い終わった燃料：止めても熱を出し続ける}
\leavevmode\label{ph:原-16:sec2}軽水炉の燃料は、ウラン235 の割合を3〜5\,\% に高めた（濃縮した）ウランです（原-15 ⑤）。燃料は原子炉の中で何年か使われ、そのあいだにウラン235 が分裂して、たくさんの\textbf{分裂片}（核分裂生成物）ができます。熱出力 \qty{3}{GW} の原子炉は1日に約 \qty{3}{kg}、1年に約1.2トンのウラン235 を分裂させるので（原-15 例題2）、燃料の中には、毎年それと同じくらいの質量の分裂片がたまっていきます。また、燃料の大部分をしめるウラン238 の一部は、中性子を吸ってプルトニウム239 に変わります。

分裂片の多くは放射性で、β線やγ線を出しながらこわれていきます。そのため、取り出した燃料は、強い放射線を出し、熱（崩壊熱）も出し続けます。この燃料を、\textbf{使用済燃料}といいます。\bkidx{しようずみねんりょう}{使用済燃料}

\phboxhead{崩壊熱は、ゆっくりとしか減らない}
原-15 ⑪の崩壊熱の目安の式を、3年運転した燃料に使ってみましょう。運転した時間は $T_0=3\times\qty{3.16e7}{s}$ です。止めてから1年後と10年後の崩壊熱の割合は
\pheq{\frac{P}{P_0}\approx0.066\left[\,t^{-0.2}-(t+T_0)^{-0.2}\,\right]=\begin{cases}0.051\,\%&（1年後）\\0.0067\,\%&（10年後）\end{cases}}{単位なし}
\label{ph:原-16:decayheat}です（例題1）。熱出力 \qty{3}{GW} の原子炉1基分の燃料なら、1年後に約 \qty{1.5}{MW}、10年後でも約 \qty{0.20}{MW} の熱を出しています。\qty{1.5}{MW} は、家庭用の電気ストーブ（約 \qty{1}{kW}）1500台分です。
\begin{bkfigure}
\begin{tikzpicture}
\begin{axis}[bkplot, xmode=log, ymode=log, width=112mm, height=58mm, xmin=6e4, xmax=4.5e9, ymin=1e-4, ymax=2,
  xtick={86400,2592000,31557600,315576000,3155760000}, xticklabels={1日,30日,1年,10年,100年},
  ytick={0.0001,0.001,0.01,0.1,1}, yticklabels={0.0001,0.001,0.01,0.1,1},
  xlabel={止めてからの時間}, ylabel={$P/P_0$（\%）}, axis lines=left, grid=both, minor grid style={draw=none}]
  \addplot[bkcurve, domain=86400:3.2e9, samples=120] {6.6*(x^(-0.2)-(x+94672800)^(-0.2))};
  \addplot[bkcurve2, dashed, domain=86400:3.2e9, samples=120] {6.6*(x^(-0.2)-(x+31557600)^(-0.2))};
  \fill[bkBlue] (axis cs:31557600,0.05056) circle (1.8pt);
  \fill[bkBlue] (axis cs:315576000,0.006735) circle (1.8pt);
  \node[font=\sffamily\footnotesize, anchor=south west] at (axis cs:3.6e7,0.056) {0.051 \%（1年後）};
  \node[font=\sffamily\footnotesize, anchor=south west] at (axis cs:3.6e8,0.0075) {0.0067 \%（10年後）};
  \node[font=\sffamily\footnotesize, anchor=north east] at (axis cs:2.2e6,0.12) {1年運転（破線）};
  \node[font=\sffamily\footnotesize, anchor=south west] at (axis cs:1.2e5,0.62) {3年運転（実線）};
\end{axis}
\end{tikzpicture}
\bkfigcaption{使い終わった燃料の崩壊熱の割合（原-15 ⑪の目安の式）。たて・横とも対数の目もり（たては10倍ずつ。3-12）。実線は3年運転、破線は1年運転した燃料。}\label{ph:原-16:fig-decayheat}
\end{bkfigure}

\begin{phyomi}[{$P/P_0\approx0.066\,[t^{-0.2}-(t+T_0)^{-0.2}]$}]
1年後から10年後までの9年で、崩壊熱は約7.5分の1にしか減りません。半減期の短い分裂片は早くなくなり、セシウム137（約30年）のような半減期の長い分裂片が残って、熱を出し続けるからです。止めた直後の6〜7\,\%（原-15 ⑪）と比べれば、1年後の熱は100分の1以下ですが、それでも冷やさずに置いておけば、熱はたまり続けます。
\end{phyomi}

\begin{phbaisuru}{運転した時間 $T_0$}
3年運転を6年運転（2倍）にすると、1年後の崩壊熱は約1.3倍、10年後は約1.8倍になります（図の破線と実線の違いも同じ向きです）。長く使った燃料ほど、半減期の長い分裂片がたくさんたまっているので、時間がたってからの熱が大きく残るのです。
\end{phbaisuru}

\begin{motto}[4]
\phboxhead{崩壊熱の式は、どこから来たか}
ある瞬間に、たくさんの核分裂がいっせいに起きたとします。そのあとの分裂片の崩壊熱は、何百種類もの半減期の指数関数（原-12）の和ですが、ウェイとウィグナーは1948年に、それが全体として、分裂から時間 $\tau$ たったときに $\tau^{-1.2}$ に比例するとみなせることを示しました（目安）。原子炉は、運転した時間 $T_0$ のあいだ、同じ割合で分裂を続けています。止める $s$ 秒前（$0\leqq s\leqq T_0$）に起きた分裂は、止めてから $t$ 秒後には、分裂から $t+s$ 秒たっています。その熱を、運転したあいだ全部足し合わせる（積分する：5-12）と
\[ \int_0^{T_0}(t+s)^{-1.2}\,\dd s=\Bigl[\,\frac{(t+s)^{-0.2}}{-0.2}\,\Bigr]_0^{T_0}=5\left[\,t^{-0.2}-(t+T_0)^{-0.2}\,\right] \]
となり、原-15 ⑪の式の形が出てきます。係数の 0.066 は、運転中の熱との比になるように合わせたものです。止めてすぐ（$t$ が $T_0$ よりずっと小さい）は、$(t+T_0)^{-0.2}$ が $t^{-0.2}$ よりずっと小さいので、崩壊熱は運転した時間にあまりよりません。ずっとあと（$t$ が $T_0$ よりずっと大きい）は、かっこの中が約 $0.2\,T_0\,t^{-1.2}$ になり、運転した時間に比例します。上の「2倍にすると」で、10年後のほうが運転した時間の影響が大きかったのは、このためです。
\end{motto}

\phboxhead{まず、水のプールで冷やす}
取り出した燃料は、まず発電所の中の水のプールに入れて、何年も（数年以上）冷やします。水は、崩壊熱を運び去るだけでなく、燃料から出るγ線や中性子をさえぎる役目もします（原-13 ①⑥）。燃料の上に何 m もの水があるので、プールのそばで作業する人が受ける線量は小さくおさえられます。プールの水は、ポンプで回して冷やし続けます。冷やせなくなるとどうなるかは、例題3で見積もります。

何年か冷やして熱が小さくなった燃料は、金属の頑丈な容器（キャスク）に入れて、空気の流れで冷やしながら保管することもできます。これを\textbf{乾式貯蔵}といいます。\bkidx{かんしきちょぞう}{乾式貯蔵}使い終わった燃料を、次の行き先が決まるまで、発電所の外でまとめてあずかる施設（\textbf{中間貯蔵施設}）も、2024年11月に青森県むつ市で動き始めました。\bkidx{ちゅうかんちょぞう}{中間貯蔵}

%% ================================================================
\bksubsection{3}{地層処分：はてな？の答え}
\leavevmode\label{ph:原-16:sec3}使い終わった燃料の行き先について、日本は、次のような方針をとっています（図\ref{ph:原-16:fig-cycle}）。
\begin{enumerate}[label={\bkcircled[bkTheme]{\arabic*}}, leftmargin=2.4\zw, itemsep=1pt, topsep=2pt]
\item 使い終わった燃料を、工場で溶かして、まだ使えるウランとプルトニウムを取り出し、もう一度燃料にする（\textbf{再処理}）。\bkidx{さいしょり}{再処理}
\item 残った、分裂片などをふくむ強い放射能の液を、ガラスの原料と混ぜて溶かし、ステンレスの容器の中で固める（\textbf{ガラス固化体}）。これが\textbf{高レベル放射性廃棄物}である。\bkidx{がらすこかたい}{ガラス固化体}\bkidx{こうれべるほうしゃせいはいきぶつ}{高レベル放射性廃棄物}
\item ガラス固化体を、30〜50年ほど冷やして熱を小さくしてから、地下深くの安定した地層に埋める（\textbf{地層処分}）。法律では、\textbf{地下 300 m 以上の深さ}の地層と決められている（特定放射性廃棄物の最終処分に関する法律 第2条）。\bkidx{ちそうしょぶん}{地層処分}
\end{enumerate}
このように、燃料をもう一度使う流れを\textbf{核燃料サイクル}といいます。\bkidx{かくねんりょうさいくる}{核燃料サイクル}再処理工場は青森県六ヶ所村で建設が進められていますが、2026年の時点では、まだ完成していません。
\begin{bkfigure}
\begin{tikzpicture}[x=1mm, y=1mm]
  \txqbox{k1}{(15,34)}{29}{ウラン鉱石}{bkGray}
  \txqbox{k2}{(48,34)}{29}{燃料をつくる\\（ウランは濃縮）}{bkOrange}
  \txqbox{k3}{(81,34)}{29}{原子炉で\\発電（原-15）}{bkRed}
  \txqbox{k4}{(114,34)}{29}{プールで\\冷やす（②）}{bkBlue}
  \txqbox{k5}{(114,17)}{29}{中間貯蔵\\空気で冷やす}{bkBlue}
  \txqbox{k6}{(81,17)}{29}{再処理}{bkPurple}
  \txqbox{k7}{(81,0)}{29}{ガラスで固める\\（ガラス固化体）}{bkPurple}
  \txqbox{k8}{(48,0)}{29}{30〜50年\\冷やす}{bkBlue}
  \txqbox{k9}{(15,0)}{29}{地層処分\\地下 300 m 以上}{bkGreen}
  \draw[txq arr] (k1) -- (k2); \draw[txq arr] (k2) -- (k3); \draw[txq arr] (k3) -- (k4);
  \draw[txq arr] (k4) -- (k5); \draw[txq arr] (k5) -- (k6); \draw[txq arr] (k6) -- (k7);
  \draw[txq arr] (k7) -- (k8); \draw[txq arr] (k8) -- (k9);
  \draw[txq back] (k6.west) -- (48,17) -- (k2.south);
  \node[txq lab, anchor=east] at (46.5,20.5) {ウラン・プルトニウム};
\end{tikzpicture}
\bkfigcaption{日本の方針の核燃料サイクル（形を見るための図）。破線は、取り出したウランとプルトニウムを、もう一度燃料にする流れ（プルトニウムは濃縮せず、ウランと混ぜて燃料にする）。}\label{ph:原-16:fig-cycle}
\end{bkfigure}

\phboxhead{なぜ、地下深くに、何万年も}
ガラス固化体の放射能は、はじめの数百年で大きく減ります。ほとんどの分裂片は半減期が短く、半減期が長いほうのセシウム137（約30年）でも、300年で約1000分の1になるからです（例題2）。けれども、使い終わった燃料には、プルトニウム239（半減期 約2万4千年）のように、半減期のとても長い物質もふくまれます。再処理でウランとプルトニウムの大部分を取り出しても、ガラス固化体には、半減期の長い物質が少し残ります。そのため、高レベル放射性廃棄物は、\textbf{数万年以上}のあいだ、人の暮らす所から離しておく必要があると考えられています（原子力発電環境整備機構の資料）。

図\ref{ph:原-16:fig-jikan}は、この時間の長さを、対数の目もりで並べたものです。どの物質でも、半減期から「1000分の1になる時間」までは、目もりのちょうど1つ分（10倍）です。
\begin{bkfigure}
\begin{tikzpicture}[x=1mm, y=1mm]
  \draw[bkInk, line width=0.8pt, -{Stealth[length=2.2mm]}] (0,0) -- (130,0);
  \foreach \x/\t in {6.12/1日,42/1年,56/10年,70/100年,84/1000年,98/1万年,112/10万年,126/100万年} {
    \draw[bkInk, line width=0.6pt] (\x,0) -- (\x,1.6); \node[txq lab, anchor=north] at (\x,-1) {\t};}
  % ヨウ素131
  \draw[bkGray!60, line width=0.4pt, densely dotted] (0,30) -- (126,30);
  \draw[bkRed!80!black, line width=2pt] (18.77,30) -- (32.77,30);
  \filldraw[fill=white, draw=bkRed!80!black, line width=0.8pt] (18.77,30) circle (1.4);
  \fill[bkRed!80!black] (31.57,28.8) rectangle (33.97,31.2);
  \node[txq lab, anchor=south] at (25.8,32) {ヨウ素131：8日→80日};
  % セシウム137
  \draw[bkGray!60, line width=0.4pt, densely dotted] (0,20) -- (126,20);
  \draw[bkOrange!80!black, line width=2pt] (62.68,20) -- (76.68,20);
  \filldraw[fill=white, draw=bkOrange!80!black, line width=0.8pt] (62.68,20) circle (1.4);
  \fill[bkOrange!80!black] (75.48,18.8) rectangle (77.88,21.2);
  \node[txq lab, anchor=south] at (69.7,22) {セシウム137：30年→300年};
  % プルトニウム239
  \draw[bkGray!60, line width=0.4pt, densely dotted] (0,10) -- (126,10);
  \draw[bkPurple!80!black, line width=2pt] (103.35,10) -- (117.35,10);
  \filldraw[fill=white, draw=bkPurple!80!black, line width=0.8pt] (103.35,10) circle (1.4);
  \fill[bkPurple!80!black] (116.15,8.8) rectangle (118.55,11.2);
  \node[txq lab, anchor=south east] at (126,12) {プルトニウム239：2.4万年→24万年};
  % ガラス固化体を冷やす
  \fill[bkBlue!25] (62.68,36) rectangle (65.79,40);
  \node[txq lab, anchor=west] at (67,38) {ガラス固化体を冷やす 30〜50年};
  \node[txq lab, anchor=west, align=left] at (0,45) {○ 半減期　■ 1000分の1（半減期の10倍）};
\end{tikzpicture}
\bkfigcaption{時間の物差し（対数の目もり）。○から■までは、どの物質でも目もり1つ分。}\label{ph:原-16:fig-jikan}
\end{bkfigure}

\begin{column}{自然の原子炉と、地層の時間}
原-12 ⑫で見たように、アフリカのガボンのオクロでは、約20億年前に、地下のウラン鉱床で、自然に核分裂の連鎖反応が起きていました。そこでできた分裂片やプルトニウムが、その後の長い年月のあいだに、地層の中でどれだけ動いたかが調べられています。このような「自然の実験」は、地層が長い時間をかけて放射性物質をどう閉じ込めるかを知る手がかりとして、地層処分の研究でも調べられてきました。何万年という時間は、人の歴史に比べればとても長い時間です。でも、20億年前の天然の原子炉の跡を今も調べられるように、地層にとっては、それほど長い時間ではありません。
\end{column}

\phboxhead{何重にもとじこめる：多重バリア}
地層処分では、放射性物質を、人がつくった壁と自然の壁の両方で、何重にもとじこめます（図\ref{ph:原-16:fig-barrier}）。放射性物質は、ガラスの中に固められています。ガラスは、金属の厚い容器（オーバーパック）に入れられ、その周りを、水を通しにくい粘土（緩衝材）で囲みます。その外側が、地下深くの岩盤です。深い地下では、地下水の流れがとてもゆっくりで、酸素も少ないので、金属がさびにくいと考えられています。もし、長い年月のあいだに放射性物質が少しずつ地下水に溶け出しても、岩盤の中を移動するあいだに、さらに減っていくように設計します。このように、いくつもの壁を組み合わせる考え方を、\textbf{多重バリア}といいます。\bkidx{たじゅうばりあ}{多重バリア}
\begin{bkfigure}
\begin{tikzpicture}[x=1mm, y=1mm]
  % 左：地表と深さ
  \fill[bkGreen!15] (0,44) rectangle (40,46);
  \draw[bkInk, line width=0.8pt] (0,46) -- (40,46);
  \fill[bkOrange!12] (0,0) rectangle (40,44);
  \draw[bkInk, line width=0.6pt] (18,46) -- (18,8) (22,46) -- (22,8);
  \fill[bkGray!30] (6,4) rectangle (34,8);
  \draw[bkInk, line width=0.6pt] (6,4) rectangle (34,8);
  \draw[txq dim] (2,46) -- (2,6);
  \node[txq lab, anchor=west, align=left] at (3.5,26) {300 m\\以上};
  \node[txq lab, anchor=north] at (20,3) {処分する坑道};
  \node[txq lab, anchor=south] at (20,47) {地表};
  % 右：まわりを囲む壁
  \fill[bkOrange!12] (52,0) rectangle (128,46);
  \node[txq lab, anchor=south west] at (53,1) {岩盤（自然の壁）};
  \filldraw[fill=bkGray!25, draw=bkInk, line width=0.6pt, rounded corners=3mm] (62,6) rectangle (118,40);
  \filldraw[fill=bkBlue!25, draw=bkInk, line width=0.8pt, rounded corners=2mm] (72,12) rectangle (108,34);
  \filldraw[fill=bkPurple!25, draw=bkInk, line width=0.6pt, rounded corners=1mm] (80,17) rectangle (100,29);
  \node[txq lab, align=center] at (90,23) {ガラス\\固化体};
  \node[txq lab, anchor=south] at (90,34.5) {金属の容器};
  \node[txq lab, anchor=south] at (90,40.5) {粘土（緩衝材）};
  \draw[bkGray, line width=0.5pt, densely dashed] (34,8) -- (62,6) (34,4) -- (62,6);
\end{tikzpicture}
\bkfigcaption{地層処分の多重バリア（形を見るための図。大きさの比は正しくない）。ガラス・金属の容器・粘土が人がつくった壁、岩盤が自然の壁。}\label{ph:原-16:fig-barrier}
\end{bkfigure}

\begin{senmonsho}[地層処分の安全評価の書き方]
\begin{bkhonbun}
処分場閉鎖後の長期安全性は，人工バリアと天然バリアから成る多重バリアシステムの閉じ込め性能に基づき，地下水移行シナリオを中心とした安全評価によって確認する．評価期間は数万年以上に及ぶため，地質環境の長期変遷と評価の不確実性を考慮する．
\end{bkhonbun}
\bkyaku{「処分場閉鎖後」は、廃棄物を埋め終えて坑道をふさいだあと。「人工バリア」はガラス・金属の容器・粘土、「天然バリア」は岩盤（図\ref{ph:原-16:fig-barrier}）。「地下水移行シナリオ」は、放射性物質が地下水に溶けて少しずつ運ばれる、という筋書きで、その道すじで線量を計算する。「不確実性」は、何万年先の予測がもつ、わからなさの幅のこと。}
\end{senmonsho}

\phboxhead{\bkhatenaref の答え}
使い終わった燃料は、分裂片の崩壊熱と強い放射線を出し続けるので、すぐに捨てることも、燃やすこともできません。燃やしても、原子核は変わらず、放射能は減らないからです（放射能を減らせるのは、時間だけです）。まず水のプールで何年も（数年以上）冷やし、そのあと、空気で冷やす容器で保管します。日本の方針では、再処理でウランとプルトニウムを取り出し、残りの高レベル放射性廃棄物をガラスで固めて、30〜50年ほど冷やしたあと、地下 300 m 以上の深さの地層に埋めます。セシウム137 などは数百年で大きく減りますが、半減期の長い物質があるので、数万年以上、人の暮らす所から離しておく必要があるのです。

その処分をする場所は、まだ決まっていません。法律では、処分する場所を、3つの段階の調査で選ぶことになっています。資料で地震・火山などの記録を調べる\textbf{文献調査}、地上からボーリングなどで地層を調べる\textbf{概要調査}、地下に施設をつくって調べる\textbf{精密調査}です（同じ法律 第6条〜第8条）。2020年から、いくつかの地域で最初の段階の文献調査が行われていますが、2026年の時点で、処分する場所は、まだ決まっていません。

\begin{phmada}[高レベル放射性廃棄物の処分地は？]
何万年も先まで、地層や地下水がどう変わるかを、どこまで確かに予測できるかについては、地質・材料・地下水の研究が続いています。海外では、フィンランドやスウェーデンが、使い終わった燃料を処分する場所を決め、地下の処分場の建設を進めています。日本でどこに処分するかは、地質の調査の結果だけでなく、その地域に暮らす人たちとの話し合いを通じて決めていくことになっていて、物理だけでは決まらない問題です。
\end{phmada}

%% ================================================================
\bksubsection{4}{安全のしくみ：止める・冷やす・閉じ込める}
\leavevmode\label{ph:原-16:sec4}原子炉の安全の考え方は、よく「止める・冷やす・閉じ込める」の3つで説明されます。
\begin{itemize}[label={\color{bkTheme}●}, leftmargin=1.6\zw, itemsep=1pt, topsep=1pt]
\item \textbf{止める}：制御棒を入れて、連鎖反応を止める。中性子の約 0.65\,\%（$\beta=0.0065$）がおくれて出ることと、温度が上がると増倍率が自然に下がるように作ることで、出力の変化はゆっくりになり、止める操作が間に合う（原-15 ⑥⑨⑩）。
\item \textbf{冷やす}：止めたあとも出る崩壊熱を、水で運び去り続ける（原-15 ⑪、この章の②）。
\item \textbf{閉じ込める}：放射性物質を、燃料・燃料の被覆管・原子炉圧力容器・原子炉格納容器・原子炉建屋と、何重もの壁の中に閉じ込める（図\ref{ph:原-16:fig-kabe}）。
\end{itemize}
\begin{bkfigure}
\begin{tikzpicture}[x=1mm, y=1mm]
  \filldraw[fill=bkGray!8, draw=bkInk, line width=0.8pt] (0,0) rectangle (78,46);
  \filldraw[fill=bkGreen!12, draw=bkGreen!60!black, line width=1pt, rounded corners=4mm] (6,3) rectangle (72,40);
  \filldraw[fill=bkBlue!12, draw=bkBlue!70!black, line width=1pt, rounded corners=3mm] (14,6) rectangle (64,33);
  \filldraw[fill=bkOrange!15, draw=bkOrange!70!black, line width=0.8pt] (24,9) rectangle (54,27);
  \foreach \x in {28,33,38,43,48} {\filldraw[fill=bkRed!45, draw=bkInk, line width=0.4pt] (\x,11) rectangle ++(2.4,12);}
  \node[txq lab, anchor=south] at (39,46.5) {原子炉建屋};
  \node[txq lab, anchor=south] at (39,40.5) {原子炉格納容器};
  \node[txq lab, anchor=south] at (39,33.5) {原子炉圧力容器};
  \node[txq lab, anchor=south] at (39,27.6) {燃料（被覆管の中）};
  \node[txq lab, anchor=west, align=left, text width=44mm] at (84,36) {\textbf{止める}\\制御棒・遅発中性子（原-15）};
  \node[txq lab, anchor=west, align=left, text width=44mm] at (84,23) {\textbf{冷やす}\\崩壊熱を水で運ぶ（原-15 ⑪）};
  \node[txq lab, anchor=west, align=left, text width=44mm] at (84,10) {\textbf{閉じ込める}\\何重もの壁（左の図）};
\end{tikzpicture}
\bkfigcaption{止める・冷やす・閉じ込める（形を見るための図）。燃料の周りを、被覆管・圧力容器・格納容器・建屋が何重にも囲んでいる。}\label{ph:原-16:fig-kabe}
\end{bkfigure}
そして、それぞれの働きについて、1つの設備が働かなくなっても次の設備で防げるように、備えを何段にも重ねておきます。これを\textbf{深層防護}（多重の防護）といいます。\bkidx{しんそうぼうご}{深層防護}

⑤で見るように、2011年の福島第一原子力発電所の事故では、「止める」はできましたが、地震と津波で電源をほとんどすべて失い、「冷やす」働きが長い時間失われました。その結果、燃料が溶け、「閉じ込める」壁の一部もこわれて、放射性物質が外へ出ました。

\phboxhead{事故のあとの規制}
事故のあと、原子力の安全を見る国の仕組みが見直されました。2012年9月、原子力を利用する側の役所から独立した\textbf{原子力規制委員会}が発足しました。2013年7月には、新しい規制の基準が施行され、想定する地震や津波を見直すことや、燃料が溶けるような重大な事故に備えた対策（電源と水を確保する設備など）を、電力会社に求めるようになりました。原子力発電所は、この基準に合うかどうかの審査を受けています。\bkidx{げんしりょくきせいいいんかい}{原子力規制委員会}

%% ================================================================
\bksubsection{5}{福島第一原子力発電所の事故：何が起きたか}
\leavevmode\label{ph:原-16:sec5}ここからは、2011年に起きた東京電力福島第一原子力発電所の事故について、環境省の基礎資料などの公的な資料に書かれている事実をもとに見ていきます。この本を読む人の中には、事故のときに避難した人や、家族や知り合いが避難した人、事故の影響のあった地域で今も暮らしている人がいると思います。読んでいて、つらいことを思い出したときは、無理をせずに読むのを休んでください。わからないことや心配なことは、家族や先生など、身近な人に話してみてください。

2011年3月11日午後2時46分、三陸沖を震源とするマグニチュード 9.0 の東北地方太平洋沖地震が起きました。地震と津波は、東日本の広い範囲に大きな被害をもたらしました。そのとき運転中だった福島第一原子力発電所の1〜3号機は、自動で止まりました（4〜6号機は、定期検査のため止まっていました）。そのあとの経過は、表\ref{ph:原-16:tab-keika}のとおりです。

原子炉の中で起きたことの物理は、原-15 ⑪で見たとおりです。炉心に水を入れられない時間が、1号機で約14時間、2号機・3号機で約6時間半続いたとみられています。崩壊熱で燃料が溶け（炉心溶融）、燃料の被覆管の金属ジルコニウムと水蒸気の反応でできた水素が建屋にたまって、爆発しました。高い温度と圧力で格納容器の閉じ込める働きが弱まり、すき間などから、放射性物質が外へ出ました。健康や環境への影響でおもに問題になったのは、ヨウ素131（半減期 約8日）、セシウム134（約2年）、セシウム137（約30年）、ストロンチウム90 です。

事故の大きさは、国際原子力事象評価尺度（INES）で、いちばん深刻なレベル7とされています。1986年の旧ソ連（今のウクライナ）のチョルノービリ原子力発電所の事故と同じレベルです。ただし、国連の科学委員会（UNSCEAR）の報告書は、チョルノービリでは炉心が爆発して大気に直接さらされたのに対し、福島第一では格納容器が溶けた燃料から出た放射性物質の大半をその中にとどめたことなどから、放射線による健康への影響は、チョルノービリよりはるかに小さかったとしています（環境省の基礎資料 5.2節）。一方、避難による暮らしや健康への影響は大きかったことを、⑥で見ます。
\begin{bkcenter}
\begin{bktable}{colspec={Q[l,m,wd=30mm] X[l,m]}, width=0.97\linewidth, rowsep=1.2pt, row{1}={bg=bkTheme!10, font=\sffamily\footnotesize}, cells={font=\small}}
日時 & できごと \\
2011年3月11日 午後 & 地震で外からの電気が止まり、非常用の発電機が動いた。そのあとの津波で、発電機や配電盤が水につかり、6号機を除いて交流の電源をすべて失った。原子炉を冷やせなくなった \\
3月11日 夜 & 国が原子力緊急事態宣言を出した。発電所から 3 km 以内に避難の指示、10 km 以内に屋内にとどまる指示 \\
3月12日 & 避難の指示が 10 km に広げられた。午後に1号機の建屋で水素爆発。避難の指示は 20 km に広げられた \\
3月14日 & 3号機の建屋で水素爆発。そのあと、3号機から流れこんだとみられる水素で、4号機でも水素爆発 \\
4月22日 & 20 km 以内を、原則として立ち入りできない区域に。その外でも、1年の線量が 20 mSv に達するおそれのある地域を、計画的に避難する区域に \\
12月 & 原子炉が安定して冷えた状態（冷温停止状態）になったことが確かめられた \\
\end{bktable}
\bktabcaption{事故の経過（環境省の基礎資料 6.1節・6.2節・9.4節による）。}\label{ph:原-16:tab-keika}
\end{bkcenter}

%% ================================================================
\bksubsection{6}{地域の暮らしと健康}
\leavevmode\label{ph:原-16:sec6}事故のあと、発電所の周りの多くの人が、ふるさとを離れて避難しました。福島県から避難した人は、いちばん多いときで約16万人（2012年5月）にのぼりました。避難の指示はその後、少しずつ解除されてきましたが、2026年5月の時点でも、約2万3千人が避難を続けています（福島県の集計）。

\phboxhead{除染と、避難の指示の解除}
地面や建物についた放射性物質を減らす作業を\textbf{除染}といいます。\bkidx{じょせん}{除染}方法は、放射性物質を「取り除く」（表面の土をけずる、落ち葉を集める、建物を洗う）、「さえぎる」（土でおおう）、「遠ざける」の3つで、原-14 ⑦で学んだ外部被ばくを減らす3原則（時間・距離・遮へい）と同じ考え方です。除染で取り除いた土などは、発電所のある大熊町・双葉町につくられた施設（除去土壌の中間貯蔵施設。②の使い終わった燃料の中間貯蔵とは別のもの）に集めて保管されています。法律で、2045年3月までに、福島県の外で最終処分することになっています。

避難の指示は、2020年3月までに、線量の高い「帰還困難区域」を除くすべての区域で解除されました。帰還困難区域の中でも、町や村の中心となる区域では、2022年から2023年にかけて解除が進みました。残る区域についても、2020年代をかけて、帰りたいと望む人が帰れるように、除染などが進められています（環境省の基礎資料 9.5節）。

\phboxhead{健康を見守る調査}
福島県は、事故のときの県民の健康を長く見守るために、県民健康調査を続けています。その中には、子どもだった人の甲状腺の検査や、心の健康・生活習慣の調査がふくまれています。

国連の科学委員会（UNSCEAR）は、2019年までに発表された研究をまとめた2020年/2021年の報告書で、住民が受けた線量の推定をもとに、「放射線被ばくが直接の原因となるような将来的な健康影響は検出できそうにない」と結論づけています。甲状腺の検査で多くのがんが見つかったことについては、放射線の影響ではなく、とても感度の高い超音波の検査を多くの人に行った結果だと評価しています（環境省の基礎資料 5.1節・5.2節）。

一方で、避難の生活は、体と心に大きな負担をもたらしました。病院や介護の施設にいた人の中には、避難の途中や避難のあとに亡くなった人がいます。住む所や仕事が変わったことで、生活習慣病が増えたり、心の不調を感じたりする人もいました。UNSCEAR も、避難した人の中で心臓や血管の病気などが増えたことを、放射線ではなく、社会や生活の変化によるものとしています。放射線から身を守るための避難そのものが、健康への大きな負担になりうることは、この事故から学ばれた大切なことの1つです。いまでは、事故のときにどう身を守るかを、放射線の害と、避難や移転による負担の両方を考えて決めることが大切だとされています（環境省の基礎資料 3.8節）。

事故のあと、地域では、食品の放射性物質の検査（原-14 ⑧の基準値）や、空間線量率の測定が続けられてきました。放射線の専門家と住民が、いっしょに測った線量をもとに話し合う取り組みも行われ、住民が自分で状況を判断できるようになることが、心の負担を軽くすることにもつながると考えられています（環境省の基礎資料 3.8節）。

%% ================================================================
\bksubsection{7}{廃炉と処理水}
\leavevmode\label{ph:原-16:sec7}福島第一原子力発電所では、原子炉を解体して、最後は施設をなくすまでの作業（\textbf{廃炉}、廃止措置）が続いています。\bkidx{はいろ}{廃炉}国が定めた中長期の計画（ロードマップ）では、事故から30〜40年後にあたる2041〜2051年に、廃止措置を終えることを目標にしています。おもな作業は、プールに残った燃料の取り出し、原子炉の中で溶けて固まった燃料（\textbf{燃料デブリ}）の取り出し、放射性物質をふくむ水の対策、廃棄物の処理です。\bkidx{ねんりょうでぶり}{燃料デブリ}
\begin{bkcenter}
\begin{bktable}{colspec={Q[l,m,wd=30mm] X[l,m]}, width=0.97\linewidth, rowsep=1.2pt, row{1}={bg=bkTheme!10, font=\sffamily\footnotesize}, cells={font=\small}}
時期 & できごと \\
2011年12月 & 1〜3号機の冷温停止状態が確かめられ、すべての号機が冷温停止状態に。それ以来、安定した状態を保っている \\
2014年12月 & 4号機のプールから、すべての燃料を取り出し終えた \\
2021年2月 & 3号機のプールから、すべての燃料を取り出し終えた \\
2023年8月 & ALPS 処理水を、海へ流し始めた \\
2024年11月、2025年4月 & 2号機から、燃料デブリを試験的に取り出した \\
2026年6月 & 2号機のプールから、燃料の取り出しを始めた（全615体。2026年10月29日までに126体。2028年度中に終える目標） \\
\end{bktable}
\bktabcaption{廃炉のおもな進み方（環境省の基礎資料 6.3節と、東京電力の廃炉の進み方の資料（2026年10月29日）による）。}\label{ph:原-16:tab-hairo}
\end{bkcenter}

燃料デブリの取り出しがむずかしいのは、デブリが強い放射線を出していて人が近づけないこと、格納容器の中のどこに、どんな形で、どれだけあるかが、まだよくわかっていないことによります。遠くから操作するロボットやロボットの腕を使って、中を調べながら少しずつ進められています。

\phboxhead{ALPS 処理水：トリチウムだけが残る}
原子炉の中に残る燃料デブリを冷やすために、今も水を入れ続けています。また、雨水や地下水が建屋に流れこみます。そのため、放射性物質をふくむ水（資料では「汚染水」とよばれる）が出続けています。建屋に入る地下水を減らす工事などで、その量は、1日に約 540 m$^3$（2014年5月）から約 70 m$^3$（2024年度）まで減りました。

この水から、多核種除去設備（ALPS）などで、セシウムやストロンチウムなど62種類の放射性物質を、環境に出すときの規制の基準より少なくなるまで取り除いた水を、\textbf{ALPS 処理水}といいます。\bkidx{あるぷすしょりすい}{ALPS 処理水}ただし、\textbf{トリチウム}（三重水素）だけは取り除けません。\bkidx{とりちうむ}{トリチウム}トリチウムは水素の同位体なので、酸素と結びついて、水そのものの一部（水の分子）になっているからです。そのため、トリチウムをふくむ水だけを、ふつうの水から取り分けることは、とてもむずかしいのです。トリチウムは β線を出しますが、そのエネルギーはとても小さく、水として体に入っても約10日で半分が体の外へ出るので、線量係数はセシウム137 の約700分の1です（原-14 ⑧）。

処理した水は、発電所の中のタンクにためられています。ALPS を使い始めたころは、急いで多くの水を処理することを優先したため、タンクの水の約7割には、トリチウム以外の物質も、環境に出すときの基準をこえて残っています（2025年9月の時点。環境省の基礎資料 6.3節）。この水は、流す前に、もう一度 ALPS で処理します（図\ref{ph:原-16:fig-alps}）。
\begin{bkfigure}
\begin{tikzpicture}[x=1mm, y=1mm]
  \tikzset{txq step/.style={draw=bkBlue!70!black, fill=bkBlue!8, line width=0.8pt, rounded corners=1.4mm, minimum height=15mm,
    text width=19mm, align=center, font=\sffamily\footnotesize, inner sep=2pt}}
  \node[txq step] (a1) at (11,0) {建屋にたまる水};
  \node[txq step] (a2) at (38,0) {セシウム・\\ストロンチウムを除く};
  \node[txq step] (a3) at (65,0) {ALPS で\\62種類を除く};
  \node[txq step] (a4) at (92,0) {タンクに\\ためる};
  \node[txq step, draw=bkGreen!60!black, fill=bkGreen!10] (a5) at (119,0) {海水で\\100倍以上に\\うすめる};
  \draw[txq arr] (a1) -- (a2); \draw[txq arr] (a2) -- (a3); \draw[txq arr] (a3) -- (a4); \draw[txq arr] (a4) -- (a5);
  \draw[txq back] (a4.south) -- ++(0,-5) -| (a3.south);
  \node[txq lab, anchor=north] at (78.5,-13) {基準をこえる水は、もう一度};
  \node[txq lab, anchor=north, align=center] at (119,-10.5) {海へ\\（1500 Bq/L 未満）};
\end{tikzpicture}
\bkfigcaption{ALPS 処理水ができるまで（環境省の基礎資料 6.3節による。形を見るための図）。トリチウムは、どの段でも取り除けない。}\label{ph:原-16:fig-alps}
\end{bkfigure}

国は、2021年4月に、ALPS 処理水を海へ流す方針を決め、2023年8月24日から放出が始まりました。方針では、流す前に海水で100倍以上にうすめて、トリチウムの濃さを1 L あたり \qty{1500}{Bq} 未満にします。この値は、原子力発電所などに使われているトリチウムの規制の基準（告示濃度限度。1 L あたり \num{60000} Bq）の40分の1で、世界保健機関（WHO）が飲み水について示している値（1 L あたり \num{10000} Bq）の約7分の1です（図\ref{ph:原-16:fig-noudo}）。1年に流すトリチウムの量は、事故の前の発電所の管理の目標と同じ、22兆 Bq（$\qty{2.2e13}{Bq}$）より少なくします。
\begin{bkfigure}
\begin{tikzpicture}[x=1mm, y=1mm]
  \foreach \x/\t in {0/0.1,15/1,30/10,45/100,60/1000,75/1万,90/10万} {
    \draw[bkGray!60, line width=0.4pt] (40+\x,-1) -- (40+\x,41); \node[txq lab, anchor=north] at (40+\x,-1.5) {\t};}
  \node[txq lab, anchor=north] at (85,-6.5) {トリチウムの濃さ（Bq/L。目もりは10倍ずつ）};
  % 値の札は棒の中（方眼の線は棒の下にかくれる。白地にしない：G-8）
  \fill[bkRed!40] (40,33) rectangle (126.7,38);
  \node[txq lab, anchor=east] at (38.5,35.5) {規制の基準（告示濃度限度）};
  \node[txq lab, anchor=east] at (125.7,35.5) {\num{60000}};
  \fill[bkOrange!45] (40,25) rectangle (115,30);
  \node[txq lab, anchor=east] at (38.5,27.5) {WHO の飲み水の値};
  \node[txq lab, anchor=east] at (114,27.5) {\num{10000}};
  \fill[bkBlue!40] (40,17) rectangle (102.6,22);
  \node[txq lab, anchor=east] at (38.5,19.5) {放出する水の上限};
  \node[txq lab, anchor=east] at (101.6,19.5) {1500 未満};
  \fill[bkBlue!25] (40,9) rectangle (77.2,14);
  \node[txq lab, anchor=east] at (38.5,11.5) {放出口の近く（予測）};
  \node[txq lab, anchor=east] at (76.2,11.5) {30 ほど};
  \fill[bkGreen!35] (40,1) rectangle (55,6);
  \node[txq lab, anchor=east] at (38.5,3.5) {ふだんの周りの海};
  \node[txq lab, anchor=east] at (54,3.5) {0.1〜1};
\end{tikzpicture}
\bkfigcaption{トリチウムの濃さの比較（環境省の基礎資料 6.3節による）。横の目もりは10倍ずつの対数の目もり。「放出口の近く」は、流したあとの海の濃さの予測で、いちばん高い所の値。}\label{ph:原-16:fig-noudo}
\end{bkfigure}

\noindent 東京電力が国際的な方法で行った評価では、ALPS 処理水を流すことで人が受ける線量は、日本人が自然放射線から受ける1年 \qty{2.1}{mSv} の、約100万分の1から約7万分の1とされました（2023年2月の評価）。放出を始めてから1年間の実際の放出量と海のデータを使った評価では、約40万分の1から約10万分の1でした。国際原子力機関（IAEA）は、日本政府の求めで安全性を調べ、2023年7月の報告書で、放出の方法が国際的な安全の基準に合っていて、人と環境への放射線の影響は無視できるほどだと結論づけています（環境省の基礎資料 6.3節）。放出にあたっては、漁業をはじめ、いろいろな立場から心配の声も上がりました。国や東京電力などは、IAEA の協力も得て、海水や魚のトリチウムなどをはかり続け、その結果を公表しています。

原-14 例題9では、1 L あたり \qty{1500}{Bq} のトリチウムをふくむ水を、毎日 2 L、1年間飲み続けたとしたら、という考えの上の計算で、預託実効線量が約 \qty{0.020}{mSv} になることを見ました。例題4では、この計算を、規制の基準の濃さと比べてみます。

\begin{phmada}[燃料デブリを、すべて取り出せるか？]
燃料デブリは、1〜3号機をあわせて約880トンあると推定されていますが、2024年と2025年の2回で試験的に取り出せたのは、合わせて約 0.9 g です。原子炉の中のどこに、どんな形で、どれだけあるかは、まだ一部しかわかっていません。強い放射線の中で動くロボット、デブリを切ったりつかんだりする方法、取り出したデブリを長く安全に保管する方法など、研究と工事が続いています。取り出したあとのデブリや、廃炉で出る廃棄物を最後にどう処分するかも、まだ決まっていません。
\end{phmada}

%% ================================================================
\bksubsection{8}{放射線の利用}
\leavevmode\label{ph:原-16:sec8}原子の力は、発電のほかにも、私たちの身の回りで広く使われています。どれも、これまでの章で学んだ放射線の性質を使っています（表\ref{ph:原-16:tab-riyou}）。

\begin{phmijika}[沖縄県のウリミバエの根絶]
ウリミバエは、ウリやゴーヤーなどの実に卵を産む害虫で、かつては沖縄県の島々に広くいました。そこで、育てた虫に放射線をあてて子どもができないようにし、大量に野に放しました。野生の虫は、子どもができない虫とつがいになるので、数がしだいに減っていきます。20年余りの取り組みのすえ、1993年に、沖縄県からウリミバエを根だやしにすることができました。薬を使わずに、ねらった1種類の虫だけを減らせる方法です。
\end{phmijika}

医療で受ける放射線は、日本では1人あたり平均で年に約 \qty{2.6}{mSv} と見積もられていて、自然放射線と同じくらいの大きさです（原-14 ⑩）。医療の放射線は、検査や治療で得られる利益が、放射線による害より大きいときに使われます。放射線を「使う」ときも、「量で考える」ことが基本になっています。
\begin{bkcenter}
\begin{bktable}{colspec={Q[l,m,wd=10mm] Q[l,m,wd=58mm] X[l,m]}, width=\linewidth, rowsep=0.6pt, colsep=4pt, row{1}={bg=bkTheme!10, font=\sffamily\footnotesize}, cells={font=\small}}
分野 & 例 & 使う性質（学んだ章） \\
医療 & X線の写真・CT 検査 & 骨と筋肉で減り方が違う（原-13 ③） \\
 & がんの放射線治療（X線・陽子線） & 電離で細胞を傷つける（原-13 ⑦⑫） \\
 & 薬の目印で体の中を見る（PET など） & 体の外で γ線をとらえる（原-11 ⑪） \\
工業 & 溶接の検査、厚さ計・密度計 & 密度で減り方が変わる（原-13 ⑪） \\
 & 医療の道具を包んだまま滅菌する & γ線が包みを通りぬける（原-13） \\
農業 & 病気に強い品種をつくる & 遺伝子に変化を起こす \\
 & 不妊にした虫を放して害虫を減らす & 虫の生殖の細胞を傷つける \\
研究 & 遺跡や岩石の年代をはかる & 半減期（原-12 ⑦） \\
\end{bktable}
\bktabcaption{放射線の利用の例。}\label{ph:原-16:tab-riyou}
\end{bkcenter}

%% ================================================================
\bksubsection{9}{原子の力とどうつきあうか}
\leavevmode\label{ph:原-16:sec9}物理シリーズの『熱学』（熱-16 ⑨）では、原子力発電について「CO$_2$ が少なく、ずっと動かせますが、事故への備えと、使い終わった燃料の扱いが大きな課題です」と書きました。この巻で学んだ物理を使うと、この2つの課題の中身を、数で説明できるようになります。
\begin{itemize}[label={\color{bkTheme}●}, leftmargin=1.6\zw, itemsep=1pt, topsep=1pt]
\item \textbf{事故への備え}：連鎖反応は止められても、崩壊熱は止められません。止めた直後で運転中の6〜7\,\%、1日後でも約 0.5\,\%（原-15 ⑪）。電源と水を失わないこと、失っても冷やし続けられることが、安全の中心になります（④⑤）。
\item \textbf{使い終わった燃料}：1年後でも約 0.05\,\%の熱（②）、セシウム137 は300年で約1000分の1、半減期の長い物質があるので数万年以上の隔離（③）。処分する場所は、まだ決まっていません。
\end{itemize}

物理は、これらの数を、正しく出すことができます。けれども、原子力発電をどれだけ使うか、処分の場所をどう決めるか、事故の影響を受けた地域をどう支えていくかは、物理の数だけでは決まりません。安全をどこまで確かめれば十分と考えるか、二酸化炭素や燃料の輸入、費用とのかね合いをどう考えるか、将来の世代に何を残すか、といった、人の価値の判断がかかわるからです。熱-16 で見たように、どれを重く見るかは、社会の選択です（表\ref{ph:原-16:tab-wakeru}）。
\begin{bkcenter}
\begin{bktable}{colspec={Q[l,m,wd=17mm] X[l,m] X[l,m]}, width=\linewidth, rowsep=0.6pt, colsep=4pt, row{1}={bg=bkTheme!10, font=\sffamily\footnotesize}, cells={font=\small}}
課題 & 物理（数）でわかること & 社会が選ぶこと \\
事故の備え & 崩壊熱の大きさと減り方（②・例題3） & どこまでの事故を考え、どれだけ備えるか \\
使った燃料 & 熱と放射能が減る速さ、さえぎる厚さ（②③） & 再処理をするか、どこに処分するか \\
放射線 & Bq・Sv の値、自然や医療と比べた大きさ（原-14） & どの値を基準にし、どう管理するか \\
エネルギー & 燃料の量・二酸化炭素・効率（熱-16） & どの発電を、どれだけ組み合わせるか \\
\end{bktable}
\bktabcaption{物理でわかることと、社会が選ぶことを分ける。}\label{ph:原-16:tab-wakeru}
\end{bkcenter}

この巻を学んだみなさんは、「放射線」「原子力」ということばを聞いたとき、「何 Bq か、何 Sv か」「半減期はどれくらいか」「何を、どの基準で数えた値か」「だれが、いつ出した資料か」を確かめることができます。数を確かめたうえで、いろいろな立場の人の考えを聞き、自分の考えをもつこと。それが、原子の力とつきあうための、物理の力です。



\begin{kotoba}\small
\bkword{使用済燃料}{原子炉で使い終わった燃料。分裂片をふくみ、強い放射線と崩壊熱を出す。}
\bkword{ガラス固化体}{再処理で残った強い放射能の液を、ガラスで固めたもの。}
\bkword{地層処分}{高レベル放射性廃棄物を、地下 300 m 以上の深さの安定した地層に埋めること。}
\bkword{多重バリア}{ガラス・容器・粘土・岩盤の何重もの壁で閉じ込める考え方。}
\bkword{深層防護}{1つの備えが働かなくても次の備えで防げるように、何段にも備えること。}
\bkword{燃料デブリ}{事故で溶けて、周りの物といっしょに固まった燃料。}
\bkword{ALPS 処理水}{トリチウム以外の放射性物質を、基準より少なくなるまで除いた水。}
\end{kotoba}

\Needspace{8\baselineskip}
\begin{phsaki}{原子力工学・放射線の利用・保健物理}
原子炉の設計と安全、再処理、廃棄物の処分、廃炉は\textbf{原子力工学}で、放射線を医療・工業・農業に使う技術は\textbf{放射線工学}や医学物理学で、人と環境を放射線から守る方法は\textbf{保健物理学}で学びます。地層処分には、地質学や材料の研究も欠かせません。エネルギーの選び方は、エネルギー工学やエネルギー政策の分野です。
\end{phsaki}

%% ================================================================
\Needspace{12\baselineskip}\bksection{reidai}{見本の計算}

\begin{phreidai}[使い終わった燃料の崩壊熱]\label{ph:原-16:ex1}
熱出力 \qty{3.0}{GW} の原子炉で3年運転した燃料を取り出しました。止めてから1年後と10年後の崩壊熱を、原-15 ⑪の目安の式で求めましょう。1年は $\qty{3.16e7}{s}$（365.25日）とします。
\phbamen
使い終わった燃料を、プールで何年冷やすかを考える場面です（②）。
\phmushi
燃料は3年間、同じ出力で運転したとします。目安の式は、精密な値と2倍ほどずれることもある式です（原-15 ⑪）。
\phzu
$P_0=\qty{3.0}{GW}$、$T_0=3\times\qty{3.16e7}{s}=\qty{9.47e7}{s}$、$t=\qty{3.16e7}{s}$（1年）・$\qty{3.16e8}{s}$（10年）。
\phriyu
止めたあとの熱は、分裂片の崩壊熱なので、原-15 ⑪の式を使います。
\phtateru
\[ P=P_0\times0.066\left[\,t^{-0.2}-(t+T_0)^{-0.2}\,\right] \]
\phtoku
1年後：$t^{-0.2}=0.03164$、$(t+T_0)^{-0.2}=(\qty{1.26e8}{})^{-0.2}=0.02398$ なので
\[ \frac{P}{P_0}=0.066\times(0.03164-0.02398)\approx\num{5.1e-4}=0.051\,\% \]
\[ P\approx\num{3.0e9}\times\num{5.1e-4}\approx\qty{1.5e6}{W}=\qty{1.5}{MW} \]
10年後：$t^{-0.2}=0.01996$、$(t+T_0)^{-0.2}=(\qty{4.10e8}{})^{-0.2}=0.01894$ なので
\[ \frac{P}{P_0}=0.066\times(0.01996-0.01894)\approx\num{6.7e-5}=0.0067\,\%,\quad P\approx\qty{0.20}{MW} \]
\phtashikame
単位：割合に W をかけて W。\\
大きさ：1年後から10年後で約7.5分の1。止めた直後（約6.4\,\%）と比べると、1年後は約130分の1です。\\
極端な場合：$t$ が $T_0$ よりずっと大きいと、かっこの中の2つの項が近くなり、崩壊熱は0に近づきます。でも、100年後でも約 0.0005\,\%（約 \qty{15}{kW}）が残ります。
\kotae 1年後 約0.051\,\%（約\qty{1.5}{MW}）、10年後 約0.0067\,\%（約\qty{0.20}{MW}）
\end{phreidai}

\begin{phreidai}[1000分の1になるまでの時間]\label{ph:原-16:ex2}
セシウム137（半減期 約30年）と、プルトニウム239（半減期 約2万4千年）が、それぞれはじめの1000分の1になるまでの時間を求めましょう。
\phbamen
使い終わった燃料や高レベル放射性廃棄物を、どれくらいの時間、人から離しておく必要があるかの桁を知りたい場面です（③）。
\phmushi
ほかの物質に変わって増える分（親から子ができる分：原-12 ⑧）は考えず、その物質がこわれて減ることだけを考えます。
\phzu
半減期 $T_{1/2}$、時間 $t$、残る割合 $N/N_0$。
\phriyu
半減期ごとに半分になるので、$N/N_0=(1/2)^{t/T_{1/2}}$（原-12 ③、3-11④）。
\phtateru
\[ \Bigl(\frac12\Bigr)^{t/T_{1/2}}=\frac1{1000}\quad\text{から}\quad \frac{t}{T_{1/2}}=\frac{\log1000}{\log2}=\frac{3}{0.301}\approx9.97 \]
\phtoku
半減期の約10回分です（$(1/2)^{10}=1/1024$）。
\[ \text{セシウム137}：\ 30\times9.97\approx300\ \text{年} \]
\[ \text{プルトニウム239}：\ \num{24110}\times9.97\approx\num{240000}\ \text{年} \]
\phtashikame
単位：年×（単位なし）＝年。\\
大きさ：ヨウ素131（約8日）なら約80日で、原-12 例題2 と同じ数です。半減期が長いほど、同じ割合に減るのに長い時間がかかります。\\
極端な場合：100万分の1になるのは、半減期の約20回分（セシウム137 で約600年）です。
\kotae セシウム137 約300年、プルトニウム239 約24万年
\end{phreidai}

\begin{phreidai}[冷やせなくなったプール]\label{ph:原-16:ex3}
例題1 の燃料（1年後の崩壊熱 \qty{1.5}{MW}）が、\qty{1000}{t} の水のプールに入っています。プールを冷やす設備が止まり、水も足せなくなったとします。水の温度が30\,℃から100\,℃になるまでの時間と、そのあと1日に蒸発する水の量を求めましょう。水の比熱を \qty{4.18}{kJ/(kg.K)}、100\,℃の水の蒸発熱を \qty{2257}{kJ/kg} とします。
\phbamen
冷やす働きを失ったとき、どれくらいの時間の余裕があるかを見積もる場面です（④）。
\phmushi
崩壊熱は、このあいだ一定とします（1年後の崩壊熱は、2日のあいだに 0.4\,\% ほどしか変わりません）。プールの壁や空気へにげる熱、燃料そのものの温まりは無視します。水はよく混ざって、全体が同じ温度とします。
\phzu
熱 $P=\qty{1.5e6}{W}$、水の質量 $m=\qty{1.0e6}{kg}$、比熱 $c$、蒸発熱 $L$、温度の上がり $\Delta\theta=\qty{70}{K}$（$T$ は運転した時間・半減期に使うので、温度は $\theta$ と書きます）。
\phriyu
温めるのに要る熱は $mc\,\Delta\theta$（熱-1 ③）。熱が1秒に $P$ ずつ入るので、時間は「要る熱÷$P$」です。蒸発は、1秒に $P/L$ の水を水蒸気にします。
\phtateru
\[ t_{100}=\frac{mc\,\Delta\theta}{P},\qquad \text{1日に蒸発する水}=\frac{P\times\qty{86400}{s}}{L} \]
\phtoku
\[ t_{100}=\frac{\num{1.0e6}\times\num{4.18e3}\times70}{\num{1.5e6}}\approx\qty{1.9e5}{s}\approx\qty{54}{h}\ (\text{約2日}) \]
\[ \frac{\num{1.5e6}\times86400}{\num{2.257e6}}\approx\qty{5.7e4}{kg}=\qty{57}{t} \]
\phtashikame
単位：kg×J/(kg·K)×K÷W＝J÷(J/s)＝s。\\
大きさ：約2日で100\,℃になり、そのあと1日に約 \qty{57}{t} ずつ水が減ります。数日の余裕はありますが、冷やすか水を足すかしなければ、いつかは燃料が水から出てしまいます。プールにも、冷やす設備と水を足す手段を何重にも用意しておく理由です（④）。\\
極端な場合：止めて1日後の燃料（崩壊熱は約 \qty{15}{MW}）なら、同じプールが約5時間で100\,℃になります。
\kotae 約2日（約\qty{54}{h}）で100\,℃、そのあと1日に約\qty{57}{t}
\end{phreidai}

\begin{phreidai}[処理水のトリチウムと規制の基準]\label{ph:原-16:ex4}
原-14 例題9 と同じく、毎日 2 L、1年（365日）飲み続けたと考えて、トリチウムの濃さが（1）1 L あたり \qty{1500}{Bq}（放出する水の上限）、（2）1 L あたり \num{60000} Bq（告示濃度限度）のときの預託実効線量を求めましょう。線量係数は \qty{0.000018}{\micro Sv}/Bq（大人）です。
\phbamen
規制の基準の値が、どんな線量を目安に決められているかを、計算で確かめる場面です（⑦）。実際に処理水を飲むわけではありません。
\phmushi
飲む水以外からのトリチウムは数えません。大人の線量係数だけで計算します。
\phzu
濃さ $C$〔Bq/L〕、1日 2 L、365日、係数 0.000018 \si{\micro Sv}/Bq、預託実効線量 $E$。
\phriyu
取りこんだ Bq の合計に線量係数をかけます（原-14 ⑧）。
\phtateru
\[ E=(C\times2\times365)\times0.000018 \]
\phtoku
（1）\[ E=(1500\times730)\times\num{1.8e-5}=\num{1.10e6}\times\num{1.8e-5}\approx\qty{20}{\micro Sv} \]
（2）\[ E=(60000\times730)\times\num{1.8e-5}=\num{4.38e7}\times\num{1.8e-5}\approx\qty{790}{\micro Sv} \]
$\qty{20}{\micro Sv}=\qty{0.020}{mSv}$、$\qty{790}{\micro Sv}\approx\qty{0.79}{mSv}$ です。
\phtashikame
単位：Bq×\si{\micro Sv}/Bq＝\si{\micro Sv}。\\
大きさ：（2）は約 \qty{0.8}{mSv} で、告示濃度限度が「その水を毎日約 2 L 飲み続けたとき、線量が平均で1年に \qty{1}{mSv}（公衆の線量の目安：原-14 ⑫）に達する濃さ」として決められていることと、桁が合います（子どもの係数を使うなど、詳しい計算は少し違います）。（1）は（2）の40分の1で、自然放射線 \qty{2.1}{mSv} の約100分の1です。\\
本物では：処理水は海でさらにうすまるので、人が受ける線量は、これよりずっと小さいと評価されています（⑦）。
\kotae （1）約\qty{0.020}{mSv}（2）約\qty{0.79}{mSv}
\end{phreidai}

%% ================================================================
\bksection{tsumazuki}{まちがえやすいところをチェック}
\begin{tsumazuki}{半減期の2倍たてば}
\bkng{半減期の2倍の時間がたてば、放射性物質は全部なくなる。}
\bkok{半減期の2倍で $\tfrac14$、10倍で約 $\tfrac1{1000}$。0にはならず、同じ割合ずつ減り続ける。}
半減期の長い物質があると、何万年も先まで考える必要があります（③）。
\end{tsumazuki}
\begin{tsumazuki}{崩壊熱の減り方}
\bkng{崩壊熱も、決まった時間ごとに半分になる。}
\bkok{崩壊熱は半減期の違う何百種類もの分裂片の熱の合計で、時間がたつほど減り方がゆっくりになる。}
1年後から10年後までの9年で、約7.5分の1にしか減りません（②）。
\end{tsumazuki}
\begin{tsumazuki}{Bq が同じなら}
\bkng{放射能（Bq）が同じなら、体が受ける線量（Sv）も同じだ。}
\bkok{線量は、取りこんだ Bq に、核種ごとの線量係数をかけて決まる。}
同じ 1500 Bq を飲みこんでも、トリチウム（水）なら約 \qty{0.027}{\micro Sv}、セシウム137 なら約 \qty{20}{\micro Sv} で、約700倍違います（原-14 ⑧）。「何 Bq か」だけでなく、「何の Bq か」を確かめます。
\end{tsumazuki}

%% ================================================================
\Needspace{12\baselineskip}\bksection{phgear}{工学ではこう使う}
\begin{phkougaku}[使い終わった燃料を、空気で冷やす：乾式貯蔵]
プールで何年か冷やして崩壊熱が小さくなった燃料は（②）、空気の流れだけで冷やせます。そこで、金属の厚い容器（キャスク）に入れてふたを閉じ、建物の中に並べて保管します。これが乾式貯蔵です。

キャスクは、いくつもの役目を1つでこなすように設計されます。厚い金属は γ線をさえぎり、中性子をさえぎる材料（水素の多い樹脂など）も組みこまれます（原-13 ⑥）。中の燃料が臨界にならないように、燃料の間に中性子を吸う材料を入れます（原-15 ④）。容器の表面からは、崩壊熱が空気へにげていきます。空気は、温められると軽くなって上へ流れ、下から冷たい空気が入ってくる（自然の対流）ので、ポンプや電源がなくても冷やし続けられます。④⑤で見た「電源を失う」ことへの備えにもなる方法です。

技術者は、崩壊熱（W）と、表面から空気へ熱が伝わる速さ（熱-11・熱-12）を比べて燃料の温度を確かめ、外の線量率（原-14）から壁の厚さを決めます。

\begin{phketa}[原子の力とつきあうときの数]
\phketaline{核分裂1回／炭素が燃える1回}{約 \qty{200}{MeV}／約 \qty{4.1}{eV}（約5千万倍）}
\phketaline{崩壊熱（3年運転、1年後／10年後）}{約 0.051\,\%／約 0.0067\,\%}
\phketaline{セシウム137 が1000分の1になる時間}{約300年（半減期の10倍）}
\phketaline{地層処分の深さ}{地下 300 m 以上（法律）}
\phketaline{処理水のトリチウムの濃さの上限}{1 L あたり \qty{1500}{Bq} 未満}
\phketaline{日本の自然放射線}{1年に約 \qty{2.1}{mSv}}
\phboxhead{答えが桁で合っているか}
崩壊熱は1日後 0.5\,\%・1年後 0.05\,\% ほど、半減期は「10回で約1000分の1」で桁をつかみます。\si{\micro Sv} と mSv（1000倍）の取り違えにも気をつけます。
\end{phketa}
\end{phkougaku}
\phstation{原子力工学,放射線工学,保健物理学,エネルギー工学}{使用済燃料、崩壊熱、乾式貯蔵、核燃料サイクル、再処理、ガラス固化体、地層処分、多重バリア、深層防護、廃止措置、燃料デブリ、ALPS 処理水、告示濃度限度、放射線利用}


\begin{furikaeri}
\begin{checklist}
\item 原子の力の使いみち（発電・医療・工業・農業・研究）と、核分裂の発見から原子炉・原子爆弾・原子力発電までの歴史を説明できる。\bkfurref{ph:原-16:sec1}
\item 使い終わった燃料が熱と放射線を出し続けるわけを、崩壊熱と半減期で説明し、冷やす期間と、人から離しておく期間の桁を見積もれる。\bkfurref{ph:原-16:ex2}
\item 福島第一原子力発電所の事故の経過と、廃炉・処理水のことを、物理の量（崩壊熱・Bq・Sv）と公的な資料の数で説明できる。\bkfurref{ph:原-16:sec5}
\item 放射線や原子力の数を、何を・どの基準で数えた値かに気をつけて読み、物理でわかることと、社会が選ぶことを分けて考えられる。\bkfurref{ph:原-16:sec9}
\end{checklist}
\end{furikaeri}
\bkfukidashi{mina}{物理は「どう使うか」までは決めてくれない。でも、崩壊熱も半減期も Sv も、数で確かめられる。数を確かめてから、いろいろな人の考えを聞いてみよう。}
